% GENERATED FILE. Edit canonical lessons, then run npm run generate:books.
% canonical-source-hash: fnv1a64:c6a3f6f3005b1b9c
% canonical-lessons: RU-C275-zatyagivatsya, RU-C275-otkladyvat, RU-C275-zaderzhivatsya, RU-C275-povtoryatsya, RU-C275-sovpadat

\chapter{Delays and coincidences}
\label{ch:ru-delays-and-coincidences}
\hlchaptermodality{275}

\begin{chapteropening}
\textbf{By the end of this chapter:} \emph{I can say that something dragged on, was put off, was held up, happened again or coincided with something else.}

Complete the last lesson of chapter 275: I can say that something dragged on, was put off, was held up, happened again or coincided with something else.
\end{chapteropening}

\section[\texorpdfstring{\ru{затягиваться} (zatyágivatsya)}{затягиваться (zatyágivatsya)}]{\ru{затягиваться} (zatyágivatsya) — to drag on}

\label{lesson:RU-C275-zatyagivatsya}



\begin{quote}
Before the new one: say the Russian for later on, then the Russian for at the very end.
\end{quote}

\subsection*{You'll want to know: \ru{затягиваться}}
\textbf{\ru{затягиваться}} (\emph{zatyágivatsya}) — \textquotedblleft{}to drag on\textquotedblright{}.

\begin{quote}
\ru{Собрание} \ru{затянулось} \ru{до} \ru{вечера}.
\end{quote}

\begin{quote}
\emph{Sobrániye zatyanúlos do véchera.}
\end{quote}

\begin{quote}
The meeting dragged on until the evening.
\end{quote}

Imperfective; perfective \emph{\ru{затянуться}} (\emph{zatyanútsya}). Of a wound it means \textquotedblleft{}to heal over\textquotedblright{}.

\begin{cousinweb}
\emph{\ru{затягиваться}} is built on \emph{\ru{за}-} and \emph{\ru{тянуть}} \textquotedblleft{}to pull\textquotedblright{}, which is inherited from Proto-Slavic \emph{tęgnǫti}, from PIE \emph{teng\textsuperscript{h}-} \textquotedblleft{}to pull\textquotedblright{} (EtymDB): to be pulled out long.
\end{cousinweb}

\begin{grammarlens}[title={What you've built}]
One more word for telling, explaining and describing at B1.
\end{grammarlens}

\subsection*{Your turn}
\begin{itemize}
\raggedright
  \item \emph{Say it:} \emph{zatyágivatsya}
  \item \emph{Say it:} \emph{zatyágivatsya}, once more
  \item \emph{Say it:} a sentence of your own with \emph{zatyágivatsya}
  \item \emph{Recall:} say the Russian for later on, then the Russian for at the very end, then say \emph{zatyágivatsya} again
\end{itemize}

\begin{tcolorbox}[breakable,colback=teal!4,colframe=teal!35!black,title={Before you move on}]
What does \ru{затягиваться} mean? (\textquotedblleft{}To drag on\textquotedblright{}.) Say it once more.
\end{tcolorbox}
\section[\texorpdfstring{\ru{откладывать} (otkládyvat)}{откладывать (otkládyvat)}]{\ru{откладывать} (otkládyvat) — to put off, to postpone; to save up}

\label{lesson:RU-C275-otkladyvat}



\begin{quote}
Before the new one: say the Russian for at the very end, then the Russian for to drag on.
\end{quote}

\subsection*{You'll want to know: \ru{откладывать}}
\textbf{\ru{откладывать}} (\emph{otkládyvat}) — \textquotedblleft{}to put off, to postpone; to save up\textquotedblright{}.

\begin{quote}
\ru{Они} \ru{откладывали} \ru{поездку} \ru{три} \ru{раза}.
\end{quote}

\begin{quote}
\emph{Oní otkládyvali poyézdku tri ráza.}
\end{quote}

\begin{quote}
They put off the trip three times.
\end{quote}

Imperfective; perfective \emph{\ru{отложить}} (\emph{otlozhít}). Also \textquotedblleft{}to put money aside\textquotedblright{}: \emph{\ru{откладывать} \ru{деньги}}.

\begin{cousinweb}
\emph{\ru{откладывать}} is built on \emph{\ru{от}-} \textquotedblleft{}away\textquotedblright{} and the stem \emph{\ru{клад}-} of \emph{\ru{класть}} \textquotedblleft{}to put\textquotedblright{}, an inherited Slavic verb (Proto-Slavic \emph{klasti}): to put to one side. EtymDB lists no parent for the compound.
\end{cousinweb}

\begin{grammarlens}[title={What you've built}]
One more word for telling, explaining and describing at B1.
\end{grammarlens}

\subsection*{Your turn}
\begin{itemize}
\raggedright
  \item \emph{Say it:} \emph{otkládyvat}
  \item \emph{Say it:} \emph{otkládyvat}, once more
  \item \emph{Say it:} a sentence of your own with \emph{otkládyvat}
  \item \emph{Recall:} say the Russian for at the very end, then the Russian for to drag on, then say \emph{otkládyvat} again
\end{itemize}

\begin{tcolorbox}[breakable,colback=teal!4,colframe=teal!35!black,title={Before you move on}]
What does \ru{откладывать} mean? (\textquotedblleft{}To put off, to postpone; to save up\textquotedblright{}.) Say it once more.
\end{tcolorbox}
\section[\texorpdfstring{\ru{задерживаться} (zadérzhivatsya)}{задерживаться (zadérzhivatsya)}]{\ru{задерживаться} (zadérzhivatsya) — to be delayed, to stay late}

\label{lesson:RU-C275-zaderzhivatsya}



\begin{quote}
Before the new one: say the Russian for to drag on, then the Russian for to put off.
\end{quote}

\subsection*{You'll want to know: \ru{задерживаться}}
\textbf{\ru{задерживаться}} (\emph{zadérzhivatsya}) — \textquotedblleft{}to be delayed, to stay late\textquotedblright{}.

\begin{quote}
\ru{Папа} \ru{сегодня} \ru{задерживается} \ru{на} \ru{работе}.
\end{quote}

\begin{quote}
\emph{Pápa sevódnya zadérzhivayetsya na rabóte.}
\end{quote}

\begin{quote}
Dad is staying late at work today.
\end{quote}

Imperfective; the perfective is \emph{\ru{задержаться}} (\emph{zaderzhátsya}): \emph{\ru{поезд} \ru{задержался} \ru{на} \ru{час}} — the train was an hour late.

\begin{cousinweb}
\emph{\ru{задерживаться}} is built on \emph{\ru{за}-} and \emph{\ru{держать}} \textquotedblleft{}to hold\textquotedblright{}, which is inherited from Proto-Slavic \emph{dĭržati} \textquotedblleft{}to hold\textquotedblright{}, from PIE \emph{d\textsuperscript{h}er-} (EtymDB): to be held back.
\end{cousinweb}

\begin{grammarlens}[title={What you've built}]
One more word for telling, explaining and describing at B1.
\end{grammarlens}

\subsection*{Your turn}
\begin{itemize}
\raggedright
  \item \emph{Say it:} \emph{zadérzhivatsya}
  \item \emph{Say it:} \emph{zadérzhivatsya}, once more
  \item \emph{Say it:} a sentence of your own with \emph{zadérzhivatsya}
  \item \emph{Recall:} say the Russian for to drag on, then the Russian for to put off, then say \emph{zadérzhivatsya} again
\end{itemize}

\begin{tcolorbox}[breakable,colback=teal!4,colframe=teal!35!black,title={Before you move on}]
What does \ru{задерживаться} mean? (\textquotedblleft{}To be delayed, to stay late\textquotedblright{}.) Say it once more.
\end{tcolorbox}
\section[\texorpdfstring{\ru{повторяться} (povtoryátsya)}{повторяться (povtoryátsya)}]{\ru{повторяться} (povtoryátsya) — to happen again, to recur}

\label{lesson:RU-C275-povtoryatsya}



\begin{quote}
Before the new one: say the Russian for to put off, then the Russian for to be delayed.
\end{quote}

\subsection*{You'll want to know: \ru{повторяться}}
\textbf{\ru{повторяться}} (\emph{povtoryátsya}) — \textquotedblleft{}to happen again, to recur\textquotedblright{}.

\begin{quote}
\ru{Такая} \ru{история} \ru{больше} \ru{не} \ru{повторится}.
\end{quote}

\begin{quote}
\emph{Takáya istóriya bólshe ne povtorítsya.}
\end{quote}

\begin{quote}
A story like this will not happen again.
\end{quote}

Imperfective; perfective \emph{\ru{повториться}} (\emph{povtorítsya}). \emph{\ru{Повторять}} is to repeat something; with \emph{-\ru{ся}} the event itself comes back.

\begin{cousinweb}
\emph{\ru{повторяться}} is built on \emph{\ru{по}-} and \emph{\ru{второй}} \textquotedblleft{}second\textquotedblright{}, which is inherited from Proto-Slavic \emph{vŭtorŭ} \textquotedblleft{}second\textquotedblright{} (EtymDB): to happen a second time.
\end{cousinweb}

\begin{grammarlens}[title={What you've built}]
One more word for telling, explaining and describing at B1.
\end{grammarlens}

\subsection*{Your turn}
\begin{itemize}
\raggedright
  \item \emph{Say it:} \emph{povtoryátsya}
  \item \emph{Say it:} \emph{povtoryátsya}, once more
  \item \emph{Say it:} a sentence of your own with \emph{povtoryátsya}
  \item \emph{Recall:} say the Russian for to put off, then the Russian for to be delayed, then say \emph{povtoryátsya} again
\end{itemize}

\begin{tcolorbox}[breakable,colback=teal!4,colframe=teal!35!black,title={Before you move on}]
What does \ru{повторяться} mean? (\textquotedblleft{}To happen again, to recur\textquotedblright{}.) Say it once more.
\end{tcolorbox}
\section[\texorpdfstring{\ru{совпадать} (sovpadát)}{совпадать (sovpadát)}]{\ru{совпадать} (sovpadát) — to coincide, to match}

\label{lesson:RU-C275-sovpadat}



\begin{quote}
Before the new one: say the Russian for to be delayed, then the Russian for to happen again.
\end{quote}

\subsection*{You'll want to know: \ru{совпадать}}
\textbf{\ru{совпадать}} (\emph{sovpadát}) — \textquotedblleft{}to coincide, to match\textquotedblright{}.

\begin{quote}
\ru{Наш} \ru{отпуск} \ru{совпал} \ru{с} \ru{их} \ru{отпуском}.
\end{quote}

\begin{quote}
\emph{Nash ótpusk sovpál s ikh ótpuskom.}
\end{quote}

\begin{quote}
Our holiday coincided with theirs.
\end{quote}

Imperfective; perfective \emph{\ru{совпасть}} (\emph{sovpást}), past \emph{\ru{совпал}}. It takes \emph{\ru{с}} + instrumental.

\begin{cousinweb}
\emph{\ru{совпадать}} is built on \emph{\ru{со}-} \textquotedblleft{}together\textquotedblright{} and \emph{\ru{падать}} \textquotedblleft{}to fall\textquotedblright{}, which is inherited from Proto-Slavic \emph{padati} (EtymDB): to fall together.
\end{cousinweb}

\begin{grammarlens}[title={What you've built}]
One more word for telling, explaining and describing at B1.
\end{grammarlens}

\subsection*{Your turn}
\begin{itemize}
\raggedright
  \item \emph{Say it:} \emph{sovpadát}
  \item \emph{Say it:} \emph{sovpadát}, once more
  \item \emph{Say it:} a sentence of your own with \emph{sovpadát}
  \item \emph{Recall:} say the Russian for to be delayed, then the Russian for to happen again, then say \emph{sovpadát} again
\end{itemize}

\begin{tcolorbox}[breakable,colback=teal!4,colframe=teal!35!black,title={Before you move on}]
What does \ru{совпадать} mean? (\textquotedblleft{}To coincide, to match\textquotedblright{}.) Say it once more.
\end{tcolorbox}
